% Zumdahl Chapter 9 -- Covalent Bonding: Orbitals.
% ELABORATED notes. Only 9.1 is on the CED; 9.2-9.5 (molecular orbital
% theory) is a clearly marked enrichment block. 3 blocks total.
\documentclass{apchem-notes}

\apcunitword{Chapter}
\apcunit{9}
\apcunittitle{Covalent Bonding: Orbitals}
\apcdoctitle{Guided Notes}
\apcsubtitle{Zumdahl \S9.1 (CED) $+$ \S9.2--9.3 (enrichment) \textbullet\ PDF pp.~454--490 \textbullet\ 3 blocks}

\begin{document}
\apctitleblock

\begin{apcnote}[Read this before you open the chapter --- it is mostly off-syllabus]
This is the \textbf{thinnest chapter in the book for AP purposes}. Only
\textbf{\S9.1} is examinable --- hybridization and $\sigma$/$\pi$ bonding,
which are CED 2.7.

\textbf{\S9.2--9.5 develop molecular orbital theory} --- MO diagrams, bond
order from filled orbitals, homonuclear and heteronuclear diatomics. None of
it is on the current CED. Rather than pretend otherwise, this chapter is
built as two examinable blocks plus \textbf{one clearly marked enrichment
block}. Work Blocks 1--2 properly; treat Block 3 as optional reading.

One more scope note: the CED assesses only \textbf{sp, sp$^2$, and sp$^3$}.
Zumdahl's \S9.1 also covers dsp$^3$ and d$^2$sp$^3$ for expanded octets ---
and he himself notes that this model ``does not give an accurate picture''
and that newer work does not use $d$ orbitals at all. You still assign
\emph{shapes} to \ce{PCl5} and \ce{SF6} with VSEPR; you are simply not asked
for their hybridization.
\end{apcnote}

% =====================================================================
\blocklesson{Hybridization and the Localized Electron Model}{Zumdahl \S9.1}

\begin{objectives}
  \item Explain why hybrid orbitals are proposed at all.
  \item Assign sp, sp$^2$, or sp$^3$ from the electron-domain count.
\end{objectives}

\reading{Zumdahl \S9.1}{455--466}

\begin{warmup}[3]
  \item Electron domains around C in \ce{CH4}: \answerblank[1.6cm]{4}
  \item Geometry of \ce{BF3}: \answerblank[3.4cm]{trigonal planar}
  \item Does a double bond count as one domain or two?
        \answerblank[2cm]{one}
\end{warmup}

\segment{INSTRUCTION A}{25}{The problem hybridization solves}

\notesectionz{Why plain atomic orbitals are not enough}{9.1}
\practice{1}

Carbon's ground-state configuration is $[\ce{He}]\,2s^2\,2p^2$ --- one
filled $s$ orbital and two half-filled $p$ orbitals, pointing at
90$^\circ$ to each other. Taken literally, that predicts carbon should form
\answerblank[1.4cm]{two} bonds at \answerblank[1.6cm]{90$^\circ$}.

Methane is observed to have \answerblank[1.4cm]{four} identical C--H bonds
at \answerblank[2cm]{109.5$^\circ$}. The atomic-orbital picture is simply
wrong about the molecule.

\begin{apcnote}
Zumdahl's resolution, and the key idea of the whole chapter: \emph{start
with the molecule, not the atoms}. A molecule's tendency to minimize its
energy matters more than preserving the orbital arrangement the free atom
happened to have. So we assume the atom \term{hybridizes} --- its $2s$ and
three $2p$ orbitals combine into four identical orbitals pointing at the
corners of a tetrahedron.

Hybridization is a \emph{model built to fit observed geometry}, not a
physical event that happens to atoms. Nothing ``hybridizes'' in time.
\end{apcnote}

\notesub{The three hybridizations you are responsible for}

\renewcommand{\arraystretch}{1.45}
\noindent\begin{tabular}{@{}cp{2.6cm}p{2.9cm}p{3.4cm}p{2.6cm}@{}}
\toprule
\textbf{Domains} & \textbf{Hybrid} & \textbf{Built from} &
\textbf{Geometry} & \textbf{Leftover $p$}\\
\midrule
2 & \answerblank[1.4cm]{sp} & one $s$, one $p$ &
  \answerblank[1.8cm]{linear} & \answerblank[1cm]{2}\\
3 & \answerblank[1.4cm]{sp$^2$} & one $s$, two $p$ &
  \answerblank[3.2cm]{trigonal planar} & \answerblank[1cm]{1}\\
4 & \answerblank[1.4cm]{sp$^3$} & one $s$, three $p$ &
  \answerblank[2.5cm]{tetrahedral} & \answerblank[1cm]{0}\\
\bottomrule
\end{tabular}

The count is the whole trick: \textbf{the number of hybrid orbitals always
equals the number of atomic orbitals combined}, and equals the number of
electron \answerblank[1.6cm]{domains}.

\segment{GUIDED PRACTICE}{15}{Zumdahl's three-step procedure}

Always in this order:
\begin{enumerate}[leftmargin=*,itemsep=0.15em]
  \item Draw the \answerblank[3.3cm]{Lewis structure}.
  \item Determine the electron-pair arrangement using
        \answerblank[1.6cm]{VSEPR}.
  \item Assign the \answerblank[3.2cm]{hybrid orbitals} that match.
\end{enumerate}

Apply it:
\begin{enumerate}[leftmargin=*,itemsep=0.3em]
  \item C in \ce{CH4}: \answerblank[4.5cm]{4 domains $\to$ sp$^3$}
  \item B in \ce{BF3}: \answerblank[4.5cm]{3 domains $\to$ sp$^2$}
  \item C in \ce{CO2}: \answerblank[4.5cm]{2 domains $\to$ sp}
  \item N in \ce{NH3}: \answerblank[5.0cm]{4 domains (1 LP) $\to$ sp$^3$}
  \item O in \ce{H2O}: \answerblank[5.0cm]{4 domains (2 LP) $\to$ sp$^3$}
\end{enumerate}

\begin{apctrap}
Count \emph{domains}, not bonds. Water's oxygen has only two bonds but two
lone pairs as well --- four domains, so sp$^3$. Students who count bonds
report sp for water and get both the hybridization and the bent geometry
wrong.
\end{apctrap}

\segment{INSTRUCTION B}{20}{The model's limits}

\notesectionz{Where the simple picture strains}{9.1}
\practice{6}

For \ce{PCl5}, five electron pairs require a trigonal bipyramidal
arrangement, and Zumdahl introduces \term{dsp$^3$} hybridization to supply
five orbitals. But he immediately adds a caution worth taking seriously:

\begin{apcnote}
Zumdahl's own words are that this model ``does not give an accurate picture
of the actual bonding,'' and that recent research shows a better model
\emph{does not use $d$ orbitals at all}. He keeps it only for convenience.

This is a good general lesson: the simple models in an introductory course
are chosen for usefulness, not truth, and a well-taught course tells you
which is which. It is also why the CED stops at sp$^3$ --- you assign
\ce{PCl5} and \ce{SF6} their \emph{shapes} with VSEPR and leave their
hybridization alone.
\end{apcnote}

\segment{APPLICATION}{20}{Assign hybridization in real molecules}

\begin{enumerate}[leftmargin=*,itemsep=0.4em]
  \item In methanol, \ce{CH3OH}, give the hybridization of both the carbon
        and the oxygen, with domain counts.
        \answerlines[2]{C: four bonds $\to$ 4 domains $\to$ sp$^3$.
        O: two bonds plus two lone pairs $\to$ 4 domains $\to$ sp$^3$.}
  \item In formaldehyde, \ce{H2CO}, give the carbon's hybridization and
        predict the H--C--H angle. \answerlines[2]{Three domains (two C--H
        bonds and one C=O) $\to$ sp$^2$, trigonal planar, angle close to
        120$^\circ$.}
  \item \ce{SO2} and \ce{SO3} both have sulfur as the central atom, yet only
        one is bent. Give each hybridization and explain.
        \answerlines[3]{\ce{SO3}: 3 domains, no lone pair $\to$ sp$^2$,
        trigonal planar. \ce{SO2}: 3 domains including a lone pair $\to$
        also sp$^2$, but the lone pair occupies one position, leaving a bent
        \emph{molecular} shape.}
\end{enumerate}

\begin{exitticket}
Both \ce{CH4} and \ce{H2O} have sp$^3$ hybridized central atoms, yet their
bond angles differ. Explain. \answerlines[2]{Both have four domains, so both
are sp$^3$. Water's two lone pairs repel more strongly than bonding pairs,
compressing its angle from 109.5$^\circ$ to about 104.5$^\circ$.}
\end{exitticket}

% =====================================================================
\blocklesson{Sigma and Pi Bonds}{Zumdahl \S9.1}

\begin{objectives}
  \item Distinguish $\sigma$ from $\pi$ bonds by how orbitals overlap.
  \item Count $\sigma$ and $\pi$ bonds, and explain restricted rotation.
\end{objectives}

\reading{Zumdahl \S9.1}{455--466}

\begin{warmup}[3]
  \item Hybridization of C in \ce{C2H6}: \answerblank[2cm]{sp$^3$}
  \item How many hybrid orbitals does sp$^2$ produce?
        \answerblank[1.6cm]{3}
  \item Leftover unhybridized $p$ orbitals on an sp carbon:
        \answerblank[1.4cm]{2}
\end{warmup}

\segment{INSTRUCTION A}{25}{Two ways for orbitals to overlap}

\notesectionz{Sigma and pi}{9.1}
\practice{1}

\renewcommand{\arraystretch}{1.4}
\noindent\begin{tabular}{@{}p{2.4cm}p{4.6cm}p{6.0cm}@{}}
\toprule
 & \textbf{$\sigma$ (sigma) bond} & \textbf{$\pi$ (pi) bond}\\
\midrule
Overlap & \answerblank[2.4cm]{head-on}, along the internuclear axis &
  \answerblank[2.6cm]{side-by-side}, above and below the axis\\
Built from & hybrid orbitals & leftover \answerblank[3.2cm]{unhybridized
  $p$} orbitals\\
Electron density & concentrated \answerblank[3.9cm]{between the nuclei} &
  in two lobes off the axis\\
Rotation & \answerblank[1.8cm]{free} & \answerblank[2cm]{blocked}\\
\bottomrule
\end{tabular}

\begin{center}
\begin{tikzpicture}[font=\sffamily\footnotesize, scale=0.95]
  % sigma
  \fill[apcgray!40] (0,0) circle (0.42);
  \fill[apcgray!40] (1.5,0) circle (0.42);
  \fill[apcgray!70] (0.75,0) ellipse (0.5 and 0.3);
  \node[circle,fill=black,inner sep=1.6pt] at (0,0) {};
  \node[circle,fill=black,inner sep=1.6pt] at (1.5,0) {};
  \node at (0.75,-0.95) {$\sigma$: overlap ON the axis};
  % pi
  \begin{scope}[xshift=5.4cm]
    \node[circle,fill=black,inner sep=1.6pt] at (0,0) {};
    \node[circle,fill=black,inner sep=1.6pt] at (1.5,0) {};
    \draw[thick] (0,0) -- (1.5,0);
    \fill[apcgray!55] (0.75,0.55) ellipse (0.95 and 0.34);
    \fill[apcgray!55] (0.75,-0.55) ellipse (0.95 and 0.34);
    \node at (0.75,-1.25) {$\pi$: lobes ABOVE and BELOW};
  \end{scope}
\end{tikzpicture}
\end{center}

\notesub{Counting rule}

\begin{itemize}[leftmargin=*,itemsep=0.2em]
  \item Single bond $=$ \answerblank[1.6cm]{1 $\sigma$}
  \item Double bond $=$ \answerblank[2.6cm]{1 $\sigma$ + 1 $\pi$}
  \item Triple bond $=$ \answerblank[2.6cm]{1 $\sigma$ + 2 $\pi$}
\end{itemize}

\begin{apcnote}
The \emph{first} bond between two atoms is always $\sigma$; every additional
bond is $\pi$. That is because head-on overlap is more effective than
side-by-side, so it happens first --- and it is also why $\pi$ bonds are
generally weaker and more reactive.
\end{apcnote}

\segment{GUIDED PRACTICE}{15}{Count them}

\begin{enumerate}[leftmargin=*,itemsep=0.35em]
  \item \ce{C2H6} (ethane): \answerblank[4.5cm]{7 $\sigma$, 0 $\pi$}
  \item \ce{C2H4} (ethene): \answerblank[4.5cm]{5 $\sigma$, 1 $\pi$}
  \item \ce{C2H2} (ethyne): \answerblank[4.5cm]{3 $\sigma$, 2 $\pi$}
  \item \ce{HCN}: \answerblank[4.5cm]{2 $\sigma$, 2 $\pi$}
  \item \ce{N2}: \answerblank[4.5cm]{1 $\sigma$, 2 $\pi$}
\end{enumerate}

\segment{INSTRUCTION B}{20}{Why double bonds cannot rotate}

\notesectionz{Restricted rotation}{9.1}
\practice{6}

A $\sigma$ bond is cylindrically symmetric about the internuclear axis, so
the two ends can spin freely --- rotation costs essentially nothing.

A $\pi$ bond is different. Its two lobes must stay
\answerblank[2.4cm]{parallel} to overlap at all. Rotating one end by
90$^\circ$ would set the $p$ orbitals perpendicular, destroying the overlap
and \answerblank[3.9cm]{breaking the bond}. That takes a large amount of
energy, so at ordinary temperatures a double bond is
\answerblank[1.6cm]{rigid}.

\begin{workedexample}[Worked example: why cis and trans exist]
In \ce{C2H4Cl2} with one Cl on each carbon, the two chlorines can sit on the
same side of the double bond (\emph{cis}) or opposite sides
(\emph{trans}). Because the $\pi$ bond blocks rotation, these are
\textbf{different compounds} --- separable, with different melting points
and dipole moments.

Do the same with ethane (\ce{C2H4Cl2} with a single C--C bond) and the two
arrangements interconvert freely billions of times a second. They are not
distinguishable substances at all.
\end{workedexample}

\segment{APPLICATION}{20}{Structure and bonding together}

\begin{enumerate}[leftmargin=*,itemsep=0.4em]
  \item For \ce{H2C=CH-C#N}: give the hybridization of each carbon and the
        total $\sigma$ and $\pi$ counts.
        \answerlines[3]{\ce{CH2} carbon: 3 domains, sp$^2$. Middle carbon:
        3 domains, sp$^2$. Nitrile carbon: 2 domains, sp. Totals:
        6 $\sigma$ and 3 $\pi$.}
  \item Explain why ethyne (\ce{C2H2}) is linear while ethene
        (\ce{C2H4}) is planar but not linear.
        \answerlines[3]{Ethyne's carbons have two domains each $\to$ sp,
        180$^\circ$, so every atom lies on one line. Ethene's carbons have
        three domains $\to$ sp$^2$, 120$^\circ$, giving a flat molecule with
        bent angles at each carbon.}
\end{enumerate}

\begin{exitticket}
A molecule has 4 $\sigma$ and 2 $\pi$ bonds between its heavy atoms. What
kind of multiple bonding must it contain? \answerlines[2]{Two $\pi$ bonds
means either one triple bond (1$\sigma$ + 2$\pi$) or two separate double
bonds (each 1$\sigma$ + 1$\pi$).}
\end{exitticket}

% =====================================================================
\blocklesson{ENRICHMENT: The Molecular Orbital Model}{Zumdahl \S9.2--9.3 --- NOT on the CED}

\begin{apcnote}[This block is optional]
Nothing in this block is assessed on the AP exam. It is included because
molecular orbital theory answers one question the Lewis model cannot ---
and knowing where a model \emph{fails} is worth an hour of anyone's time.
Skip it entirely without penalty.
\end{apcnote}

\begin{objectives}
  \item Explain what a molecular orbital is and compute bond order.
  \item Describe the oxygen paramagnetism problem and its resolution.
\end{objectives}

\reading{Zumdahl \S9.2--9.3}{466--476}

\segment{INSTRUCTION A}{25}{Orbitals that belong to the molecule}

\notesectionz{Bonding and antibonding}{9.2}
\practice{1}

The localized electron model treats bonds as belonging to \emph{pairs} of
atoms. Molecular orbital theory instead combines atomic orbitals into
orbitals that span the \emph{whole molecule}.

Combining two atomic orbitals always produces two molecular orbitals:
\begin{itemize}[leftmargin=*,itemsep=0.2em]
  \item a \term{bonding} MO --- lower in energy, with electron density
        \answerblank[3.9cm]{between the nuclei};
  \item an \term{antibonding} MO ($\sigma^*$, $\pi^*$) --- higher in energy,
        with a \answerblank[1.6cm]{node} between the nuclei.
\end{itemize}

\[
  \text{Bond order} =
  \answerblank[6.5cm]{$\dfrac{(\text{bonding e}^-) - (\text{antibonding e}^-)}{2}$}
\]

A bond order of zero means \answerblank[3cm]{no bond forms} --- which is
exactly why \ce{He2} does not exist.

\segment{GUIDED PRACTICE}{15}{Bond order practice}

\begin{enumerate}[leftmargin=*,itemsep=0.3em]
  \item \ce{H2}: 2 bonding, 0 antibonding $\to$ bond order
        \answerblank[1.4cm]{1}
  \item \ce{He2}: 2 bonding, 2 antibonding $\to$ bond order
        \answerblank[1.4cm]{0}
  \item \ce{O2}: 10 bonding, 6 antibonding $\to$ bond order
        \answerblank[1.4cm]{2}
  \item \ce{N2}: 10 bonding, 4 antibonding $\to$ bond order
        \answerblank[1.4cm]{3}
\end{enumerate}

Zumdahl notes the payoff: across the period-2 diatomics, as MO bond order
rises, bond energy \answerblank[1.8cm]{increases} and bond length
\answerblank[1.8cm]{decreases} --- strong evidence the model is capturing
something real.

\segment{INSTRUCTION B}{20}{The oxygen problem}

\notesectionz{Where Lewis structures fail}{9.3}
\practice{6}

\begin{apcnote}
Draw the Lewis structure of \ce{O2} and you get \ce{O=O} with two lone pairs
on each oxygen --- \textbf{every electron paired}. That structure predicts
oxygen should be \term{diamagnetic}, weakly repelled by a magnet.

Liquid oxygen is \term{paramagnetic}. Pour it between the poles of a magnet
and it visibly sticks. It has \textbf{two unpaired electrons}.

The Lewis model cannot account for this. It is not a small discrepancy or a
measurement issue --- the model simply gets it wrong.
\end{apcnote}

The MO model resolves it cleanly. Filling oxygen's molecular orbitals places
the last two electrons in \emph{separate} degenerate $\pi^*$ orbitals, with
parallel spins by Hund's rule. Result: bond order 2 (agreeing with the
double bond) \emph{and} two \answerblank[2.2cm]{unpaired} electrons,
matching the magnetic evidence exactly.

\segment{APPLICATION}{20}{Comparing the two models}

\begin{enumerate}[leftmargin=*,itemsep=0.4em]
  \item State one thing the localized electron model does better than MO
        theory, and one thing MO does better.
        \answerlines[3]{Localized electron model: far easier to apply, and
        it predicts molecular geometry, which MO theory does not do
        naturally. MO theory: correctly predicts magnetic behaviour and
        handles delocalized electrons and fractional bond orders.}
  \item Both models agree that \ce{O2} has a bond order of 2. Why is that
        agreement reassuring rather than trivial?
        \answerlines[2]{Two independently constructed models converging on
        the same measurable quantity is evidence that the quantity is real
        rather than an artifact of either model's assumptions.}
\end{enumerate}

\begin{exitticket}
Why does \ce{He2} not exist, in MO terms? \answerlines[2]{Its two bonding
and two antibonding electrons give a bond order of zero --- there is no net
stabilization, so no bond forms.}
\end{exitticket}

\end{document}
